\section{Scaling Law y Emergence: cómo una pérdida total suave oculta cambios abruptos en las tareas}

La sección anterior descompuso el objetivo de entrenamiento en una mezcla de tareas latentes; esta sección estudia el eje de escala. La segunda y la tercera intuición de Jason no se contradicen: al aumentar la cantidad de cómputo, la cross-entropy total puede bajar de forma estable, mientras que la accuracy de un benchmark puede permanecer en cero durante mucho tiempo y luego subir de golpe. Para entender este fenómeno hay que observar al mismo tiempo la probabilidad continua, el peso de cada tarea y la regla de calificación discreta.

\subsection{Definición operativa mínima de scaling}

Este apartado explica primero el eje horizontal. En clase, el compute de entrenamiento se entiende de forma aproximada como el producto del tamaño del modelo por la cantidad de datos; el entrenamiento real también depende de la longitud de secuencia, del cómputo de las pasadas hacia adelante y hacia atrás, del aprovechamiento del hardware y de la receta de entrenamiento, pero esta simplificación basta para expresar la dirección de «invertir más cómputo efectivo». Lo importante es comparar una serie de puntos dentro del mismo paradigma de entrenamiento, no comparar de forma aislada los nombres de dos modelos.

\lecturefigure{jason-slide-008.jpg}{Intuición dos: el cómputo de entrenamiento que representan un modelo y unos datos más grandes reduce la loss de forma confiable}{Jason Wei official deck, p.~8}

Una aproximación común es
\[
C\propto N D,
\]
donde $C$ representa la cantidad de cómputo de entrenamiento, $N$ el número de parámetros sin contar los de embedding y $D$ el número de tokens de entrenamiento; la constante de proporcionalidad absorbe las operaciones hacia adelante y hacia atrás por token. Esta expresión no es un modelo de costo completo, pero le recuerda al investigador que aumentar solo los parámetros sin aumentar los datos, o aumentar solo los datos con un modelo de capacidad insuficiente, puede alejarlo de la región compute-optimal.

\lecturefigure{jason-slide-009.jpg}{La loss de un modelo de lenguaje sigue una tendencia predecible de ley de potencias respecto al cómputo}{Jason Wei official deck, p.~9}

El ajuste típico de una scaling law se escribe
\[
L(C)=L_{\infty}+A C^{-\alpha},\qquad A>0,\ \alpha>0.
\]
donde $L(C)$ es la pérdida de validación para un compute dado, $L_{\infty}$ es la cota inferior de la pérdida, irreducible o todavía no eliminada, $A$ controla la altura de la curva y $\alpha$ controla la pendiente de la mejora conforme crece el compute. En escala logarítmica, el intervalo de ley de potencias se aproxima a una recta, por lo que pueden usarse experimentos más pequeños para estimar la loss esperada de entrenamientos más grandes.

\begin{knowledgebox}{Lectura de la figura: confirmar primero los ejes, el rango y si hay saturación}
El eje horizontal abarca varios órdenes de magnitud y el eje vertical es la loss, donde menor es mejor. La evidencia visual decisiva no es que un punto vaya adelante, sino que la tendencia no se curva de forma apreciable hacia una horizontal en un rango muy amplio. Esto respalda que «dentro del intervalo observado, seguir aumentando el cómputo todavía rinde beneficios», pero no demuestra que la tendencia continúe para siempre ni que el beneficio alcance a compensar el costo de entrenamiento.
\end{knowledgebox}

\teachervoice{En 00:09:36--00:10:44, Jason recuerda que esta relación se llama ``law'' porque sigue siendo predecible a lo largo de unos siete órdenes de magnitud, no porque, como una ley física, nunca deje de cumplirse. Un informe de investigación debe dar a la vez el intervalo de ajuste y la incertidumbre de la extrapolación.}

\subsection{Por qué funciona el scaling: hay que separar la evidencia de las conjeturas}

El apartado anterior presentó una relación empírica; este discute los mecanismos posibles. Las dos explicaciones que se dan en clase se califican expresamente de hand-wavy: un modelo más grande puede memorizar más hechos de la cola larga y también puede representar funciones más complejas en una sola forward pass. Son hipótesis útiles, pero no pueden confirmarse solo a partir de la curva de loss.

\lecturefigure{jason-slide-010.jpg}{Intuición: el modelo pequeño memoriza de forma selectiva; el modelo grande da cabida a la cola larga y ejecuta cálculos más complejos}{Jason Wei official deck, p.~10}

\begin{warningbox}{Errores causales frecuentes al leer la figura}
«Un modelo más grande tiene menor loss» no basta para demostrar que «toda la reducción proviene de la memorization», ni que «la profundidad se encarga del razonamiento y el ancho de los hechos». El número de parámetros, la cobertura de los datos, la estabilidad de la optimización y el tiempo de entrenamiento cambian al mismo tiempo. Para poner a prueba un mecanismo hay que controlar variables, construir datos contrafactuales o medir directamente cómo cambian los hechos de distinta frecuencia y las tareas que exigen distinta profundidad de cálculo.
\end{warningbox}

\teachervoice{En 00:11:13--00:12:55, Jason dice dos veces que estas son solo hand-wavy answers. Conservar esta capa de incertidumbre importa más que ofrecer un mecanismo que suene completo: el hecho del curso es que la relación de scaling es robusta; «memoria de la cola larga + funciones más complejas» es una explicación candidata.}

\subsection{La loss total es una mezcla ponderada de las pérdidas por tarea}

Volvamos ahora a la perspectiva massive multi-task de la sección anterior. Si los distintos tokens del corpus requieren, respectivamente, gramática, hechos, sentimiento o matemáticas, la loss total se comporta como un promedio ponderado de las pérdidas de esas subdistribuciones. Las tareas fáciles pueden acercarse a la saturación mientras las difíciles todavía mejoran con rapidez; por eso, que la curva total sea suave no obliga a que cada componente lo sea.

\lecturefigure{jason-slide-011.jpg}{Intuición tres: la loss total baja de forma suave, pero una tarea posterior individual todavía puede mostrar un emergent pattern}{Jason Wei official deck, p.~11}

\lecturefigure{jason-slide-012.jpg}{Descomposición de la loss total en pérdidas por tarea, como grammar, world knowledge y math}{Jason Wei official deck, p.~12}

En términos conceptuales,
\[
L_{\mathrm{all}}(C)=\sum_{k=1}^{K} w_k L_k(C),
\qquad w_k\ge 0,\quad \sum_{k=1}^{K}w_k=1.
\]
donde $K$ es el número de familias de tareas latentes, $w_k$ es el peso de esa tarea entre los tokens de evaluación y $L_k(C)$ es la pérdida de la tarea cuando el compute es $C$. Aunque la pérdida $L_j$ de una tarea difícil baje rápidamente en un intervalo estrecho, si $w_j$ es muy pequeño la curva total todavía puede verse suave.

\begin{knowledgebox}{Qué es una emergent ability}
Esta clase adopta una definición conductual: un modelo pequeño obtiene en cierta tarea un resultado cercano al azar o de cero, mientras que un modelo más grande alcanza un resultado claramente mejor. Esta definición describe \textbf{la forma de la curva de evaluación}; no implica por sí misma que dentro del modelo haya surgido de repente un módulo completamente nuevo. Si la función de calificación tiene un umbral, usa exact match o exige acertar todos los pasos de un proceso de varios pasos, una mejora continua de la probabilidad también se traduce en un cambio abrupto aparente.
\end{knowledgebox}

\subsection{Las tareas fáciles se saturan y las difíciles arrancan tarde}

El apartado anterior dio la fórmula de la mezcla; este lee en la figura las curvas que la componen. La gramática puede ser ya suficientemente buena a una escala menor, de modo que seguir agrandando el modelo solo aporta mejoras mínimas; en cambio, las tareas difíciles, como las matemáticas, pueden cruzar el umbral de utilidad solo a una escala mayor. Al leer la figura hay que comparar las pendientes y la posición de la saturación, no interpretar cada línea como un módulo independiente.

\lecturefigure{jason-slide-013.jpg}{Formas de scaling distintas para la loss total, una tarea ya saturada y una tarea difícil}{Jason Wei official deck, p.~13}

Si la calificación de la tarea se obtiene a partir de una confianza continua $q(C)$ mediante un umbral $\tau$,
\[
S(C)=\mathbf{1}\{q(C)\ge \tau\},
\]
donde $S(C)$ es la calificación discreta, $q(C)$ es la confianza de éxito, que varía de forma suave con el compute, y $\tau$ es el umbral de calificación. Aunque $q(C)$ suba de forma suave, $S(C)$ salta de 0 a 1 al cruzar $\tau$. Esta es una de las fuentes de la emergence, pero no la única.

\lecturefigure{jason-slide-014.jpg}{Las tareas de BIG-Bench incluyen curvas smooth, flat, inverse, noisy y emergent-looking}{Jason Wei official deck, p.~14}

\begin{knowledgebox}{Lectura de la figura: la distribución de las tareas informa más que la dicotomía «emerge o no emerge»}
En la figura, cerca del 29\% de las tareas mejora de forma suave, el 22\% es casi plano dentro del rango observado, el 2\% muestra inverse scaling, el 13\% no tiene una relación clara con la escala y el 33\% es emergent-looking. Los porcentajes dependen del conjunto de modelos, de las tareas y de la metric; la conclusión más sólida es que las curvas por tarea adoptan muchas formas y que la tendencia única de la loss total no sustituye la medición tarea por tarea.
\end{knowledgebox}

\teachervoice{En 00:16:54--00:19:24, Jason enfatiza que, si solo se entrena hasta antes del punto de cambio abrupto, el investigador predecirá erróneamente que esa capacidad nunca aparecerá. En la sesión de Q\&A añade que la falta de puntos intermedios densos de modelos hace muy difícil determinar si ya había señal antes del cambio abrupto.}

\subsection{La emergence en el prompting y la salvedad sobre la metric}

Además de las curvas por tarea, el propio prompt también puede revelar diferencias de escala. Un modelo pequeño puede ser incapaz de extraer la regla de correspondencia a partir de unas cuantas demostraciones, mientras que uno grande puede aprovechar el contexto para resolver entradas nuevas; pero, tras observar un salto en exact-match, todavía hay que revisar la token-level probability, las metric continuas y la robustez frente al prompt, para no describir equivocadamente la geometría de la función de calificación como un mecanismo interno.

\lecturefigure{jason-slide-015.jpg}{Los ejemplos de entrada y salida en el prompt pueden activar un in-context behavior que depende de la escala}{Jason Wei official deck, p.~15}

\begin{warningbox}{«La emergence es un mirage» y «la capacidad es falsa» no son la misma afirmación}
En la sesión de Q\&A, Jason reconoce que cambiar la metric cambia la forma de la curva; aun así, sostiene que los comportamientos útiles que muestran los modelos grandes son reales. Una formulación rigurosa distingue tres niveles: si la distribución de salida del modelo cambia de forma continua, si la calificación de la evaluación cambia de golpe y si las tareas que el usuario puede resolver experimentan un cambio cualitativo. Los tres pueden cumplirse a la vez sin contradicción.
\end{warningbox}

\teachervoice{La respuesta de 00:29:53--00:30:21 no ofrece una etiqueta simple de bando; recomienda leer los artículos y juzgar por cuenta propia. Por eso estas notas no tratan la emergence como una cuestión de creencia, sino que consideran la metric sensitivity parte del diseño experimental.}

\subsection{Resumen de la sección}

La scaling law describe cómo la pérdida total mejora de forma estable con la cantidad de cómputo dentro del intervalo observado; la emergent-looking curve describe el comportamiento de una tarea específica después de pasar por la transformación de la metric. Ambas pueden surgir de la combinación de la mezcla de tareas, las diferencias de dificultad, la calificación por umbral y el muestreo disperso. El investigador debe informar a la vez la pérdida continua, la calificación por tarea, los puntos de tamaño de modelo y las barras de error, e indicar si cada explicación de un mecanismo es evidencia, hipótesis o intuición personal.
